% =============================================================================
\chapter{Introduction}
\label{ch:introduction}
% =============================================================================

\section{Background and Motivation}
\label{sec:intro:background}

According to the World Health Organization, approximately \num{466000000}
people worldwide live with disabling hearing loss, and around 7~million of
these reside in the Arab region~\cite{who2024hearing}. Sign language is the
first language of most deaf individuals; in the Arab world, this is
\gls{arsl}. Unlike spoken languages, \gls{arsl} is a visual--spatial
language with grammar in three dimensions, and is therefore not a
one-to-one mapping of spoken Arabic.

% TODO: Figure 1-1 — pie chart of hearing-impaired population by region.
% \begin{figure}[H]\centering\includegraphics[width=0.6\textwidth]{fig_1-1_hearing_loss.png}
% \caption{Estimated hearing-impaired population by region (WHO, 2024).}
% \label{fig:hearing_loss}\end{figure}

Automated \gls{slr} promises \num{24}-hour, scalable accessibility that
human interpretation cannot match in cost or availability. It is therefore
particularly valuable for education, telemedicine, customer service, and
public-service kiosks.

\section{Problem Statement}
\label{sec:intro:problem}

Given a continuous RGB video stream from a single uncalibrated camera, the
system must automatically recognise either (a)~the letter being
finger-spelled per frame, or (b)~the isolated word being signed across a
1--\SI{2}{\second} window, in either \gls{asl} or \gls{arsl}, with
end-to-end latency below \SI{200}{\milli\second}. The interface must run
in a standard web browser and on a commodity Android phone.

\section{Aim and Objectives}
\label{sec:intro:objectives}

The aim of this project is to design, train, evaluate, and deploy a
real-time bilingual \gls{asl}/\gls{arsl} recognition system covering both
letters and words.

The specific objectives are:
\begin{enumerate}
  \item Curate four datasets covering two languages and two recognition
        levels.
  \item Extract 21- and 33-landmark vectors using Google MediaPipe.
  \item Train four deep models: two MLPs (letters), one BiLSTM with
        temporal attention, and one stacked BiLSTM--LSTM model with a
        TimeDistributed spatial encoder.
  \item Quantise the four models to \gls{tflite} for on-device inference
        (each $\leq$ \SI{10}{\mega\byte}).
  \item Implement a FastAPI backend exposing both \gls{rest} and \gls{ws}
        endpoints.
  \item Implement a React web client that performs MediaPipe hand
        detection in the browser.
  \item Implement a React-Native + Expo mobile client with offline mode.
  \item Evaluate accuracy (Top-1, Top-5, macro-F1), latency, frame rate,
        and conduct a small user study.
\end{enumerate}

\section{Scope and Limitations}
\label{sec:intro:scope}

The system covers static finger-spelling, isolated-word recognition,
real-time inference, two languages (\gls{asl} and \gls{arsl}), a single
signer with a single RGB camera, and both single-handed and two-handed
signs. It does \emph{not} cover continuous sign-language translation
(planned as future work), grammar parsing, depth or sEMG sensors,
multi-signer scene analysis, or sign languages other than \gls{asl} and
\gls{arsl}.

\section{Engineering Standards and Design Considerations}
\label{sec:intro:standards}

The system follows the following standards and best practices:

\begin{itemize}
  \item HTTP/1.1 (RFC~9110) for \gls{rest} endpoints~\cite{rfc9110}.
  \item WebSocket (RFC~6455) for streaming inference~\cite{rfc6455}.
  \item HTTPS / TLS~1.3 for transport security.
  \item W3C WCAG~2.1~AA for web accessibility with ARIA roles.
  \item ISO/IEC~25010:2011 software-quality model
        (functional suitability, performance efficiency, usability,
        reliability, security, maintainability, portability).
  \item TensorFlow SavedModel and \gls{tflite} formats for model
        portability~\cite{tflite}.
  \item OpenAPI~3.1 for machine-readable API documentation.
  \item JWT (RFC~7519) for stateless authentication.
\end{itemize}

\section{Constraints}
\label{sec:intro:constraints}

The system was designed under the following constraints:

\begin{description}
  \item[Economic] Free-tier cloud (Railway 500\,h/month, Vercel hobby plan).
  \item[Hardware] Training on an NVIDIA MX150 (2\,GB VRAM) with Kaggle P100
        used for bigger runs.
  \item[Data] \gls{karsl}-502 provides only $\approx 8$ samples per class
        --- severe data scarcity for word recognition.
  \item[Time] One academic year (October to June).
  \item[Ethical] Signers' video data used only with the consent established
        by the dataset creators.
  \item[Privacy] On-device option for users who cannot upload video.
  \item[Environmental] Lightweight models reduce energy footprint.
\end{description}

\section{Research Contributions}
\label{sec:intro:contributions}

This work makes four contributions:

\begin{enumerate}
  \item[\textbf{C-1}] \textbf{Bilingual unified pipeline.} The first
        open-source \gls{asl}+\gls{arsl} letters+words system with a single
        inference backend.
  \item[\textbf{C-2}] \textbf{ArSL-Word~v2 architecture.} A TimeDistributed
        spatial encoder followed by two \gls{bilstm}s and an \gls{lstm},
        trained on a 258-D pose-and-hands feature vector, that outperforms
        a hands-only \gls{bilstm} baseline on \gls{karsl}-502.
  \item[\textbf{C-3}] \textbf{Custom \texttt{TemporalAttention} layer.} An
        additive attention re-implemented for TensorFlow~2.10 that adds
        \textless~1\,\% parameters and improves Top-5 \gls{asl}-word
        accuracy.
  \item[\textbf{C-4}] \textbf{Deployed cross-platform product.}
        Demonstrates that the trained models are usable on a phone
        (\gls{tflite}, $\approx 5$\,MB each, $\geq 25$\,\gls{fps}) and a
        browser (server inference, $\leq 150$\,ms round-trip).
\end{enumerate}

\section{Thesis Organisation}
\label{sec:intro:organisation}

\Cref{ch:literature} reviews related work; \cref{ch:methodology} presents
datasets, feature extraction, and the four model architectures;
\cref{ch:implementation} details the backend, web, mobile, and
deployment stacks; \cref{ch:results} presents accuracy, ablation, and
real-time performance results; \cref{ch:discussion} interprets the
findings; \cref{ch:conclusion} summarises contributions; and
\cref{ch:future} outlines future research directions.
